% !TEX root = /home/eb/.vscode/extensions/qftrules.programme-colle/tmp/Exercice.tex
\begin{exo}[PC][1][oral]{Indice optique}
  Grâce à deux fibres optiques, on amène des faisceaux lumineux dans deux cuves identiques remplies d'air, de longueur $\ell=\SI{20}{cm}$, et placées côte à côte. Les fibres optiques sont distantes de $a=\SI{0.5}{mm}$ en entrée des cuves. On observe ce qui se passe sur un écran placé à une distance $D=1 \mathrm{~m}$ en aval.

  \begin{quest}
    \item Décrire la figure d'interférences. Calculer l'interfrange. On donne $\lambda=\SI{0.6}{\mu m}$.
    \sol{
    Même figure que les fentes d'Young, différence de marche $\delta = n_{\rm air}\frac{ax}{D}$, interfrange $i=\frac{\lambda D}{n_{\rm air} a}$.
    }

    \item On vide une cuve et on la remplit d'un gaz d'indice $n_{2}>n_{\text {air}}$. Comment évolue la figure d'interférences? Le décalage observé est de 70 franges, calculer $n_{2}-n_{\rm air}$.
    \sol{
    Différence de marche $\delta = n_{\rm air}\frac{ax}{D} + (n_2-n_{\rm air}) \ell $, dans l'hypothèse $x\ll D$. Le centre de la figure, correspondant à $\delta = 0$, est décalé de $x_0=0$ à $x_0=-(n_2-n_{\rm air})\frac{D \ell}{a n_{\rm air}}$.
    On sait que ce décalage correspond à $N=70$ interfranges en valeur absolue, soit $x_0 = N i \Rar (n_2-n_{\rm air})\frac{D \ell}{a n_{\rm air}} = \frac{N \lambda D}{n_{\rm air} a}$ et
    $$
    n_2-n_{\rm air} = \frac{N \lambda}{\ell} = \num{0.02}.
    $$
    }

  \end{quest}
\end{exo}
%%%%%%%%%%%%%%%%%%%%%%%%%%%%%%%%%%%

%%%%%%%%%%%%%%%%%%%%%%%%%%%%%%%%%%%
\begin{exo}[PC][1][colle]{Épaisseur du cellophane}
  \Annale[Oral physique PC]{Centrale}{2024}
Pour déterminer l'épaisseur d'un film de cellophane on propose deux méthodes.
\begin{quest}
  \item	
  \textbf{1ère méthode}. Un carré du film de 40 cm de côté est placé entre deux carrés d'aluminium de mêmes dimensions. L'ensemble constitue un condensateur dont les carrés d'aluminium forment les armatures.
  
  On place en série une résistance $R=\SI{0,8}{k\Omega}$ et un générateur qui délivre une tension créneau. On mesure la tension $u(t)$ aux bornes du dipôle.
  on constate une croissance exponentielle avec un temps caractéristique de l'ordre de $\SI{0,3}{ms}$.
  On considère que les lois de Maxwell restent valables en remplaçant $\varepsilon_0$ par $\varepsilon_0 \varepsilon_r$.  Déterminer l'épaisseur $e$ du cellophane.
  \item \textbf{2ème méthode}. On dispose d'un interféromètre de Michelson réglé en lame d'air. Sur l'une des voies on place le film de cellophane. On place un spectromètre en sortie de l'interféromètre et on observe $17 \pm \num{0,5}$ cannelures entre \SI{400}{nm} et \SI{750}{nm}.
  Par ailleurs on place le film de cellophane sur un cube de verre d'indice $n=\num{1,626}$ et on constate qu'il y a une réflexion totale sur le cellophane pour une angle d'incidence $\alpha_{\text {lim }}=36^{\circ}$.
  Déterminer l'épaisseur $e$ du cellophane.
\end{quest}

\end{exo}
%%%%%%%%%%%%%%%%%%%%%%%%%%%%%%%%%%%

%%%%%%%%%%%%%%%%%%%%%%%%%%%%%%%%%%%
\begin{exo}[PC][2][oral]{Interférences}
  Deux sources cohérentes distinctes A et B sont séparées d'une distance $AB=a$ dans le plan focal objet d'une lentille convergente. Un écran est placé dans le plan focal image de la lentille, de distance focale $f'$.
  \begin{fig}[]
    \centering
    \includegraphics[width=0.4\textwidth]{lentille_interf.pdf}
  \end{fig}
  \begin{quest}
  \item Tracer les rayons partant de A et B après la lentille.
  \sol{
  \begin{fig}[]
    \centering
    \includegraphics[width=0.4\textwidth]{lentille_interf.pdf}
  \end{fig}
  }
  \item Exprimer l'angle entre l'axe optique et le rayon partant de B après la lentille.
  \sol{
  $\theta \simeq \tan \theta = \frac{AB}{OF} = \frac{a}{f'}$.
  }
  \item Déterminer le champ d'interférence (hauteur, longueur).
  \sol{
  \begin{fig}[]
    \centering
    \includegraphics[width=0.4\textwidth]{lentille_interf.pdf}
  \end{fig}
  Longueur $L=H \tan \theta = \frac{Ha}{f'}$ avec $H$ la hauteur de la lentille. Hauteur $H \frac{L-d}{L}$ à la distance $d$ derrière la lentille. Inconvénient du dispositif : champ d'interférence s'étrécit avec la distance ! Avantage : comme pour le montage de Fraunhofer (lentille en entrée), l'approximation d'écran lointain devient exacte ici.
  }
  \end{quest}
  Soit un point $M$ dans la zone d'interférence.
  \begin{quest}
  \item Déterminer l'expression du déphasage.
  \sol{
  $\Delta\vphi = \frac{ax}{f'}$ d'après le théorème de Malus.
  }
  \item En déduire l'intensité lumineuse et l'interfrange.
  \sol{
  Formule de Fresnel des interférences, $I(x) = 2I_0\left[1+\cos\left(\frac{2\pi ax}{\lambda f'}\right)\right]$, d'interfrange
  \eqbox{i = \frac{\lambda f'}{a}.}
  }
\end{quest}
  % Trois autres questions avec un fluide s'écoulant dans une partie du champ d'interférences.
\end{exo}
%%%%%%%%%%%%%%%%%%%%%%%%%%%%%%%%%%%
